\documentclass[11pt]{letter}
\usepackage{graphicx}
\usepackage{fancyhdr}
\usepackage{geometry}

\usepackage[normalem]{ulem}
\usepackage{xcolor}
\newcommand{\edit}[2]{\textcolor{purple}{\sout{#1} \textbf{#2}}}

\geometry{tmargin=1in}
\geometry{bmargin=1in}
\geometry{lmargin=1in}
\geometry{rmargin=1in}
\geometry{headsep=1in}
\setlength\topmargin{-2.5in}
%\geometry{margin=1in}
%\geometry{headheight=0.5in}
%\fancypagestyle{firstpage}{\fancyhf{}
\fancypagestyle{firstpage}{\fancyhf{}
\setlength\topmargin{-1.25in}
\fancyhead[C]{\includegraphics[width=\paperwidth,keepaspectratio=true]{ciera_header.pdf}}
\fancyfoot[C]{\includegraphics[width=\paperwidth,keepaspectratio=true]{ciera_footer.png}}
\renewcommand{\headwidth}{\paperwidth}
\fancyheadoffset{\marginparwidth}
\fancyfootoffset{\marginparwidth}}
\renewcommand{\headrulewidth}{0pt}
% Address is commented out so I can use an electronic signature below
% If address is not commented out, then change \fancypagestyle{first} to
% \fancypagestyle{empty}
%\address{Nicholas Chapman \\
%         1212 Central Street APT 1S\\
%         Evanston, IL 60201}
%\signature{Nicholas Chapman} % commented out because I use png signature file

%They want to know:
%are you a good fit to the department? -- paragraph 2 addresses this.
%are you productive? -- cv addresses this?
%is your work going to have a big impact in the future?

\begin{document}
\begin{letter}{
               Faculty Search Committee \\
               Department of Physics \\
               University of Oregon
           }

\opening{Dear Committee:}

    I am applying for the position of Assistant Professor of Physics at the University of Oregon Department of Physics.
    With this letter, please find my curriculum vitae, statement of research plans, statement of teaching, and statement on diversity.

    Astrophysical convection is ubiquitous, but is poorly understood and therefore offers exciting research avenues.
    The well-studied problem of Rayleigh-B\'{e}nard convection (RBC) lacks many of the complexities of astrophysical convection.
    For example, stars and planets have stratified atmospheres, global rotation, and long-timescale transients.
    On the other hand, simulations of astrophysical convection which include radiation hydrodynamics are often so complex that they become like observations themselves and are again difficult to gain understanding from.
    My research sits between these extremes: I try to study comprehensible physics problems which include one or some of the nuances in astrophysical research which I then work to describe in detail.
    I have found that the same fundamental scaling laws in RBC appear in density-stratified convective flows.
    I have gained insight into how to vary fundamental control parameters to separate turbulence and rotational constraint in rotating convection.
    I have extensively studied the interactions of fast and slow convective processes, paralleling the fast convection and slow evolution seen in stars.

    %Dig deep into one problem to give a flavor of the type of analyses you do.
    %Why this is an important problem? (Obs that didn't match with theory. For the first time, I have come up with a reason to make this work -- adam's paper says we're in the right ballpark. First a-priori theory / confirmed by numerical simulations that explains the latest observations.)
    % Well on the way to solving decades-long problem.
    %Can tell the story about how we came up with this: I was running simulations of something different, we noticed that this was changing. Used previous work on understanding long-term secular thermal evolution to realize that this process was happening (convection zone was growing to a certain size but it would stop). Looking at the literature, we found this work by  Zahn and Roxburgh from the 1980s-1990s that gives an analytic description of what's going on, based off averaging over the underlying equations, structure, etc. We modified their theory to take into account the effects of viscous dissipation which is still important even at large Reynolds (0th law of turbulence), and found good agreement with simulations, etc.
    %Talk about the applied math that went into this interesting solution (below)
    %Pen + paper theory (roxburgh and zahn stuff)
    %Dedalus 3D simulations that test that theory
    %Boils down the theory to a 1D parameterization that is more broadly useful to the astrophysics community
    %Has led to fruitful new follow-up collaborations.

    %Career plan: rotation, composition gradients, magnetic fields, stratification, geometry.
    My most influential work focuses on penetrative convection.
    Observations suggest that stellar structure models and 1D prescriptions of convection underestimate the size of convection zones.
    To match these observations, boundary mixing processes implemented in 1D stellar structure models have to be finely tuned and varied from one type of star to another.
    I have developed the first \emph{a-priori} theory which can explain these observations, and follow-up work by Dr.~Adam Jermyn suggests that this theory takes a step towards solving this decades-old problem.
    I discovered this process while running simulations in which I expected very little mixing at the convective boundary, and I was surprised when the convection zone advanced well beyond the expected boundary.
    Using my previous work on the long-term secular evolution of convection zones, I realized that I was witnessing a long-timescale relaxation process.
    Zahn and Roxburgh's work on penetrative convection in the 1980s and 90s provided me with analytical descriptions of this process.
    I modified their theory to account for the effects of viscous dissipation, which cannot be neglected even in the low-dissipation turbulent regime of astrophysical flows (per the zeroth law of turbulence).
    This theory agreed well with laminar and turbulent three-dimensional simulations that I ran using the Dedalus code.
    I parameterized this theory for inclusion in 1D stellar structure models, and I am currently working on an international collaboration with Dr.~Cole Johnston to implement this theoretical prescription into 1D MESA stellar evolution models.
    In my research statement, I discuss several possible directions for the development of this work.

    Throughout my academic career, I have striven not only to conduct state-of-the-art research, but also to develop my pedagogical and mentoring expertise through workshop attendance and practice.
    I now provide research mentorship to five graduate students across two institutions (Northwestern and Univ.~Colorado).
    All of these mentorship relationships have led to collaborative contributions on papers which are published or in prep (marked on my CV).
    Within my research mentoring relationships and more broadly within my academic institution, I am dedicated to providing a just, equitable and inclusive teaching and research environment.
    During my postdoctoral fellowship at CIERA, I chaired the K12 education and public outreach taskforce, and I now serve as a core member of CIERA's JEDI (Justice, Equity, Diversity, Inclusion) committee, which has organized an external sociosystemic organizational development plan, including a departmental climate survey, in the coming academic year.
    I am excited to continue my EDI work through involvement in programs like STEM CORE at the University of Oregon, and I'd love to participate in the professional development workshops offered through the Science Literacy Program to build upon my knowledge of best teaching and mentoring practices.

    I am particularly excited about this University of Oregon faculty position.
    I bring knowledge expertise in astrophysical and geophysical convective flows that are currently not represented in the Physics department, but which could be used to forge connections and collaborations with the department of earth sciences, e.g., Profs.~Dufek, Erickson, Karlstrom, and Paty.
    Beyond this, I am an expert at evolving PDEs using spectral methods and the \emph{Dedalus} pseudospectral framework.
    Using this expertise, I see many opportunities for producitve collaboration within the Physics department in areas including biophysics and biofluidics (e.g., Profs.~Parthasarathy, Paulose, and Ursell), and would be interested in expanding the collaboration network to experts in other departments at UO like Prof.~Sutherland in Biology.
    I am also extremely interested in characterizing the evolution of massive stars and the distribution of their compact object remnants (massive stars and black holes), and would be excited to work alongside Profs.~Farr and Frey in this endeavor.
%    and the potential for productive collaboration between myself and THESE PEOPLE AT YOUR DEPARTMENT on topics like LIST HERE.
%    spectral coding methods, convection, interactions between convection stable regions, rotation, magnetism, and multi-timescale processes, expertise with the open-source Dedalus code, and also multi-scale processes such as my ``entropy rain'' research.
    I value teaching and mentorship, and am an engaged member of my department community.
    For these reasons, I am an ideal candidate for this post.

    In support of my application, please feel free to contact the following individuals:

    \footnotesize
    \begin{tabular}{ll}
        \textbf{Prof.~Benjamin P.~Brown}                        &\hspace{0.25in} \textbf{Prof.~Daniel Lecoanet}                     \\
        University of Colorado, Boulder                         &\hspace{0.25in} Northwestern University    \\
        \hspace{0.2in}Dept.~Astrophysical \& Planetary Sciences &\hspace{0.45in} Dept.~Engineering Sciences \& Applied Mathematics \\
        Tel: +1 (303) 653-2371                                    &\hspace{0.45in} CIERA                          \\
        Email: bpbrown@colorado.edu                             &\hspace{0.25in} Tel: +1 (608) 335-3950 \\
                                                                &\hspace{0.25in} Email: daniel.lecoanet@northwestern.edu             \\

        \textbf{Dr.~Matteo Cantiello}                          & \\
         Flatiron Insitute                                     & \\
         \hspace{0.2in}Center for Computational Astrophysics   & \\
         Princeton University                                  & \\
         \hspace{0.2in}Dept.~Astrophysical Sciences & \\
         Tel: +1 (805) 280-8175  & \\
         Email: mcantiello@flatironinstitute.org & \\
    \end{tabular}

    \normalsize
    If there are any other questions or concerns please do not hesitate to contact me.
    Thank you for your time and consideration.

\closing{Sincerely,}
\vspace{-0.9in}
\fromsig{\includegraphics[scale=0.2]{signature.png}}\\
\fromname{Evan H. Anders}
\end{letter}

\end{document}
